\begin{abstract}
We represent hospital records from Stanford's EHRSHOT dataset as networks in which each node is a hospital stay and two stays are connected when the patients' diagnosis histories are very similar. Using standard network measures, we ask whether connected stays share outcomes for two questions: will the stay last a week or more, and will the patient be newly diagnosed with high blood pressure within a year? We compare an \emph{honest} network, built only from diagnoses recorded before the moment of prediction, with a \emph{cheating} network that also uses later diagnoses, and with a randomly rewired network. Honest networks consist of many tight groups, and connected stays share outcomes far more often than chance. Adding future diagnoses barely changes this for hospital stays, but clearly strengthens it for high blood pressure, because the later record contains the very diagnosis being predicted. A second network, linking diagnoses that tend to follow one another, recovers recognisable families of conditions such as heart failure and sepsis.
\end{abstract}
